\section{Spectral sequences}
\label{section-spectral-sequence}

\noindent
A nice discussion of spectral sequences may be found in
\cite{Eisenbud}. See also \cite{McCleary}, \cite{Lang}, etc.

\begin{definition}
\label{definition-spectral-sequence}
Let $\mathcal{A}$ be an abelian category.
\begin{enumerate}
\item A {\it spectral sequence in $\mathcal{A}$} is given by a
system $(E_r, d_r)_{r \geq 1}$ where each $E_r$ is an object
of $\mathcal{A}$, each $d_r : E_r \to E_r$ is a morphism such
that $d_r \circ d_r = 0$ and $E_{r + 1} = \Ker(d_r)/\Im(d_r)$
for $r \geq 1$.
\item A {\it morphism of spectral sequences}
$f : (E_r, d_r)_{r \geq 1} \to (E'_r, d'_r)_{r \geq 1}$ is
given by a family of morphisms $f_r : E_r \to E'_r$ such that
$f_r \circ d_r = d'_r \circ f_r$ and such that $f_{r + 1}$
is the morphism induced by $f_r$ via the identifications
$E_{r + 1} = \Ker(d_r)/\Im(d_r)$
and
$E'_{r + 1} = \Ker(d'_r)/\Im(d'_r)$.
\end{enumerate}
\end{definition}

\noindent
We will sometimes loosen this definition somewhat and allow $E_{r + 1}$
to be an object with a given isomorphism
$E_{r + 1} \to \Ker(d_r)/\Im(d_r)$.
In addition we sometimes have a system $(E_r, d_r)_{r \geq r_0}$
for some $r_0 \in \mathbf{Z}$ satisfying the properties of the definition above
for indices $\geq r_0$. We will also call this a spectral sequence since by
a simple renumbering it falls under the definition anyway.
In fact, the cases $r_0 = 0$ and $r_0 = -1$ can be found in the literature.

\medskip\noindent
Given a spectral sequence $(E_r, d_r)_{r \geq 1}$ we define
$$
0 = B_1 \subset B_2 \subset \ldots \subset B_r \subset \ldots
\subset Z_r \subset \ldots \subset Z_2 \subset Z_1 = E_1
$$
by the following simple procedure. Set $B_2 = \Im(d_1)$
and $Z_2 = \Ker(d_1)$. Then it is clear that
$d_2 : Z_2/B_2 \to Z_2/B_2$. Hence we can define $B_3$ as the unique
subobject of $E_1$ containing $B_2$ such that $B_3/B_2$ is the image
of $d_2$. Similarly we can define $Z_3$ as the unique subobject of
$E_1$ containing $B_2$ such that $Z_3/B_2$ is the kernel of $d_2$.
And so on and so forth. In particular we have
$$
E_r = Z_r/B_r
$$
for all $r \geq 1$. In case the spectral sequence starts at $r = r_0$
then we can similarly construct $B_i$, $Z_i$ as subobjects in $E_{r_0}$.
In fact, in the literature one sometimes finds the notation
$$
0 = B_r(E_r) \subset B_{r + 1}(E_r) \subset B_{r + 2}(E_r) \subset \ldots
\subset Z_{r + 2}(E_r) \subset Z_{r + 1}(E_r) \subset Z_r(E_r) = E_r
$$
to denote the filtration described above but starting with $E_r$.

\begin{definition}
\label{definition-limit-spectral-sequence}
Let $\mathcal{A}$ be an abelian category.
Let $(E_r, d_r)_{r \geq 1}$ be a spectral sequence.
\begin{enumerate}
\item If the subobjects $Z_{\infty} = \bigcap Z_r$
and $B_{\infty} = \bigcup B_r$ of $E_1$ exist then we define
the {\it limit}\footnote{This notation is not universally accepted. In some
references an additional pair of subobjects
$Z_\infty$ and $B_\infty$ of $E_1$ such that
$0 = B_1 \subset B_2 \subset \ldots \subset B_\infty \subset Z_\infty
\subset \ldots \subset Z_2 \subset Z_1 = E_1$
is part of the data comprising a spectral sequence!}
of the spectral sequence to be the object
$E_{\infty} = Z_{\infty}/B_{\infty}$.
\item We say that the spectral sequence {\it degenerates at $E_r$}
if the differentials $d_r, d_{r + 1}, \ldots$ are all zero.
\end{enumerate}
\end{definition}

\noindent
Note that if the spectral sequence degenerates at $E_r$, then
we have $E_r = E_{r + 1} = \ldots = E_{\infty}$ (and the limit
exists of course). Also, almost any abelian category we will encounter
has countable sums and intersections.

\begin{remark}[Variant]
\label{remark-allow-translation-functors}
It is often the case that the terms of a spectral sequence have
additional structure, for example a grading or a bigrading.
To accommodate this (and to get around certain technical issues)
we introduce the following notion. Let $\mathcal{A}$ be an
abelian category. Let $(T_r)_{r \geq 1}$ be a
sequence of {\it translation} or {\it shift} functors, i.e.,
$T_r : \mathcal{A} \to \mathcal{A}$ is an isomorphism of categories.
In this setting a {\it spectral sequence} is given by a system
$(E_r, d_r)_{r \geq 1}$ where each $E_r$ is an object of
$\mathcal{A}$, each $d_r : E_r \to T_rE_r$
is a morphism such that $T_rd_r \circ d_r = 0$ so that
$$
\xymatrix{
\ldots \ar[r] &
T_r^{-1}E_r \ar[r]^-{T_r^{-1}d_r} &
E_r \ar[r]^-{d_r} &
T_rE_r \ar[r]^{T_r d_r} &
T_r^2E_r \ar[r] & \ldots
}
$$
is a complex and $E_{r + 1} = \Ker(d_r)/\Im(T_r^{-1}d_r)$ for $r \geq 1$.
It is clear what a {\it morphism of spectral sequences}
means in this setting. In this setting we can still define
$$
0 = B_1 \subset B_2 \subset \ldots \subset B_r \subset \ldots
\subset Z_r \subset \ldots \subset Z_2 \subset Z_1 = E_1
$$
and $Z_\infty$ and $B_\infty$ (if they exist) as above.
\end{remark}










